% 332_Kontinuumskonstruktionen_De_ch.tex
% Dok. 332, FFGFT-Korpus, August 2026 — Kapitel-Datei
% Zwei Wege vom Diskreten zum Kontinuum: Kategorienfehler-Vermeidung
% und Kontinuumslimes im Vergleich
%
% Hinweis zur Textart: Dieses Dokument ist ein allgemein gehaltener
% Vergleichsaufsatz, kein originärer FFGFT-Beweistext. Es referiert und
% kommentiert ein externes Preprint (Krüger 2026); die epistemischen
% Marker [K]/[B]/[S]/[SETZUNG] des Korpus werden daher hier NICHT auf
% eigene FFGFT-Aussagen angewandt, sondern nur als neutrale Hervorhebung
% von Kernaussagen benutzt.

\definecolor{ffgftblue}{RGB}{30,90,160}
\definecolor{proofgreen}{RGB}{0,110,50}
\newtcolorbox{quellbox}[1][]{breakable,enhanced,boxrule=0.6pt,arc=3pt,
  fonttitle=\bfseries\small,coltitle=white,top=4pt,bottom=4pt,left=6pt,right=6pt,
  colback=blue!5,colframe=ffgftblue,colbacktitle=ffgftblue,#1}
\newtcolorbox{kernbox}[1][]{breakable,enhanced,boxrule=0.6pt,arc=3pt,
  fonttitle=\bfseries\small,coltitle=white,top=4pt,bottom=4pt,left=6pt,right=6pt,
  colback=green!5,colframe=proofgreen,colbacktitle=proofgreen,#1}

% ============================================================
\begin{quellbox}[title={Quelle}]
Diese Abhandlung referiert und kommentiert Aussagen aus:
Krüger, M. (2026). \emph{Dimension-Graded Geometry and Typed Emergence on
Aperiodic Carriers: From Cellular Degree, Holonomy, and Regge Curvature
Support to a Falsifiable Spacetime-Limit Programme in an HLV Testbed.}
Zenodo, DOI: \href{https://doi.org/10.5281/zenodo.22105698}{10.5281/zenodo.22105698}.
Alle direkten Zitate im ersten Teil stammen aus diesem Preprint bzw.\ der
begleitenden Zusammenfassung des Autors. Die Präzisierung in
Abschnitt~2.1 (,,zwei verschiedene Bedeutungen von Kontinuum``) geht auf
eine Rückmeldung von M.~Krüger (IPI, 26.~August 2026) zu einer früheren
Fassung dieses Dokuments zurück; siehe Dok.~190, Registereinträge R96, R97 und R98.
\end{quellbox}

\section*{Einleitung}

In der theoretischen Physik taucht immer wieder dieselbe methodische Aufgabe
auf: Man möchte aus einer diskreten oder kompakten Struktur — einem Gitter,
einem Simplizialkomplex, einer kompakten Mannigfaltigkeit — ein Kontinuum
gewinnen, das sich physikalisch interpretieren lässt (Zeit, Raum, ein
Eichfeld, Krümmung). Der Übergang vom Diskreten zum Kontinuierlichen ist
dabei keineswegs automatisch erlaubt. Er erfordert eine explizit
deklarierte, mathematisch kontrollierte Abbildung, die zeigt, \emph{welche}
Struktur erhalten bleibt und \emph{welche} verlorengeht. Wird diese Abbildung
nicht angegeben, entsteht ein Kategorienfehler: Objekte aus verschiedenen
mathematischen Kategorien (zum Beispiel Ketten unterschiedlichen Grades in
einem Kettenkomplex) werden stillschweigend gleichgesetzt, ohne dass die
Gleichsetzung selbst gerechtfertigt wäre.

Diese Abhandlung stellt zwei unterschiedliche Strategien gegenüber, mit
denen sich ein solcher Übergang rechtfertigen lässt: erstens eine
sorgfältige, mehrstufige Konstruktion innerhalb der diskreten
Regge-Kalkül- und homologischen Lattice-Vakuum-Tradition, die den
Kategorienfehler durch explizite Lemmata blockiert und den
Kontinuumslimes über eine Verfeinerungsfolge mit Mosco-Konvergenz
\emph{postuliert} (nicht beweist); und zweitens eine Konstruktion, bei der
das Kontinuum nicht als Grenzwert einer Folge, sondern als kanonischer
Funktionen-Pullback entlang einer festen Überlagerungsprojektion
erscheint.

% ============================================================
\section{Die diskrete Konstruktion — Kettenkomplexe, Regge-Krümmung und der bedingte hyperbolische Limes}

\subsection{Das methodische Fundament: Kettenkomplexe ohne physikalische Vorannahme}

Ausgangspunkt ist ein endlicher orientierter Kettenkomplex $K$ mit
Randabbildungen $\partial_p$, die $\partial^2=0$ erfüllen. Die dualen
Kokettenräume $C^p=\mathrm{Hom}(C_p,\mathbb{R})$ tragen Korandoperatoren
$d_p$ mit $d^2=0$. Entscheidend ist die Feststellung, dass der ,,Grad`` $p$
hier ,,nichts weiter als dieser Kettengrad`` ist — ,,keine physikalische
Bedeutung ist eingebaut``.

\begin{kernbox}[title={Kategorienfehler-Blockade}]
$C^p\cong C^q$ für $p\neq q$ ist ,,nicht wohlgetypt, solange keine explizit
deklarierte Abbildung $\Phi:C^p\to C^q$ existiert, die nachweislich die
relevanten Operationen erhält.`` Das Lemma ,,blockiert damit reflexartig
\glqq Kante = Zeit\grqq{}, \glqq Fläche = Eichfeld\grqq{}, \glqq 3-Zelle =
Gravitation\grqq{}, bevor irgendjemand das behaupten kann, ohne die
Brücke zu liefern.``
\end{kernbox}

\subsection{Eine Dimension allein trägt keine Krümmung}

Ein Standardresultat besagt: ,,In 1D ist der Riemann-Tensor ein 2-Form mit
Werten im Endomorphismenbündel — jede 2-Form auf einer 1-Mannigfaltigkeit
verschwindet trivial.`` Wichtig: ,,extrinsische Krümmung (eine Linie
gebogen in $\mathbb{R}^2$ oder $\mathbb{R}^3$) ist etwas anderes und
widerspricht dem Satz nicht`` — dies blockiert den Fehlschluss, eine
sichtbar gekrümmte Linie liefere automatisch einen 1D-Gravitations\-freiheitsgrad.

\subsection{Kanten-Transport, Flächen-Holonomie und die Grenze reiner Eichtransformationen}

Eine $U(1)$-Phasen-1-Kokette $\theta_{ij}$ auf Kanten erzeugt über $d\theta$
eine Loop-Phase $\Phi_f$ auf Flächen, invariant unter Knoten-Rephasing
$\theta_{ij}\mapsto\theta_{ij}+\alpha_j-\alpha_i$. Ist $\theta=d\alpha$
exakt, dann $d\theta=d^2\alpha=0$ — ,,reine Knotenphasendifferenzen
erzeugen also keinen nichttrivialen Fluss auf einem zusammenziehbaren
Gebiet.`` Dies ist der erste Punkt, an dem ,,eine Fläche zum ersten Mal
eine algebraisch andersartige Information tragen kann als eine Kante`` —
ob dies je zu einem physikalischen Eichfeld wird, bleibt explizit offen.

\subsection{Regge-Krümmung sitzt auf Kodimension zwei — und die Falle der lokalen Tetraedrisierung}

Im Regge-Kalkül konzentriert sich Krümmung in einer $d$-dimensionalen
stückweise-flachen Simplizialmannigfaltigkeit auf $(d-2)$-dimensionalen
,,Hinges`` — in 3D auf Kanten. Der Defektwinkel
$\varepsilon_h=2\pi-\sum_\sigma\theta_{\sigma h}$ ist ,,nur definiert, wenn
ein gültiger globaler stückweise-flacher Komplex mit kompatiblen
Simplex-Metriken existiert.``

\begin{kernbox}[title={Die entscheidende Voraussetzung}]
,,Lokale Tetraedrisierung einzelner projizierter Zellen reicht nicht, um
ein globales Regge-Krümmungsfeld zu definieren, wenn verschiedene
Tetraeder sich mit positivem Volumen überlappen oder nicht face-to-face
treffen — dann ist der Stern eines vermeintlichen Hinge gar nicht der
Stern einer gültigen Mannigfaltigkeit, und der Defektwinkel hat keine
eindeutige geometrische Interpretation mehr.``
\end{kernbox}

Überlappen sich Tetraeder mit positivem Volumen (statt die sogenannte
HFF-Bedingung — ,,face-to-face``: Schnitte verschiedener abgeschlossener
Simplizes sind leer oder eine gemeinsame Fläche — zu erfüllen), bricht die
geometrische Bedeutung des Defektwinkels zusammen, selbst wenn die
einzelnen lokalen Rechnungen fehlerfrei sind.

Diese Einschränkung betrifft den klassischen, stückweise-flachen
face-to-face-Regge-Kalkül. Es existieren verallgemeinerte Regge-Zugänge
in der Literatur — etwa Konstruktionen mit konstant gekrümmten statt
flachen Simplizes \cite{BahrDittrich2009} oder mit variabler, sich über
die Zeit ändernder Triangulierung \cite{DittrichHoehn2012} —, doch diese
lösen ein anderes Problem: Sie ersetzen die Flachheitsannahme einzelner
Simplizes bzw.\ erlauben eine wechselnde Diskretisierung, ändern aber
nichts an der Voraussetzung einer global face-to-face-konsistenten
Triangulierung und adressieren nicht den hier beschriebenen Fall sich
überlappender, nicht face-to-face treffender Simplizes.

\subsection{Algebraische Kettenstruktur vs.\ geometrische Realisierbarkeit}

Ein zuvor beobachtetes empirisches Scheitern — im Preprint als
,,007F-Ergebnis`` bezeichnet — wird zum Konstruktionsprinzip erhoben: Eine
dort berichtete numerische Untersuchung fand nach Angabe des Autors 5455
zertifizierte Fälle, in denen sich projizierte Zellen mit positivem
Volumen überlappen. Es wird zwischen dem abstrakten
Provenienz-Kettenkomplex $K_{\mathrm{abs}}$ (algebraisch gültig allein
durch Orientierung und $\partial^2=0$, unabhängig von jeder Einbettung)
und der eingebetteten Realisierung $\iota:K_{\mathrm{abs}}\to E_\parallel$
unterschieden, die die strikte HFF-Bedingung erfüllen muss.

\begin{kernbox}[title={Saubere Isolierung des Scheiterns}]
,,$\partial^2=0$ in $K_{\mathrm{abs}}$ impliziert $d^2=0$ und alle
endlichen algebraischen Hodge-Relationen — unabhängig davon, ob eine
konkrete Realisierung $\iota$ die strikte HFF-Bedingung erfüllt. Was
scheitert, ist nur die Beförderung dieser algebraischen Zellen zu einer
global eingebetteten Mannigfaltigkeit.``
\end{kernbox}

\subsection{Hodge-Positivität}

Standard-DEC-Aussage: $\langle\omega,\Delta_p\omega\rangle=\|d_p\omega\|^2+\|\delta_p\omega\|^2\geq0$,
also $\Delta_p$ positiv-semidefinit und selbstadjungiert im deklarierten
endlichen Skalarprodukt.

\subsection{Spektrale Dimension: Trivialität an beiden Enden}

Für einen endlichen positiv-semidefiniten Operator $A$ auf einem Raum der
Dimension $N=\dim\mathcal H<\infty$ mit Eigenwerten $\lambda_1,\dots,\lambda_N$
wird $P(\sigma)=N^{-1}\mathrm{Tr}\,e^{-\sigma A}$ definiert, daraus
$d_s(\sigma)=-2\sigma P'/P$.

\begin{kernbox}[title={Zwingende Konsequenz}]
$\lim_{\sigma\downarrow0}d_s(\sigma)=0$ (Taylor $P\approx1-\sigma\bar\lambda$)
und $\lim_{\sigma\to\infty}d_s(\sigma)=0$, falls $A$ einen nichttrivialen
Nullraum hat. ,,Ein positiv-dimensionales Plateau kann auf einem
endlichen Träger nur ein Zwischenskalen-Phänomen sein — beide Endpunkte
sind trivial.``
\end{kernbox}

\subsection{Die Kontinuums-Verpflichtung über Mosco-Konvergenz — Postulat, kein Beweis}

Ein fester Träger mit beschränktem Grad/Kantengewicht hat einen
beschränkten Graph-Laplace — keine intrinsische UV-Folge $\lambda\to\infty$.
Mit schrumpfender Skala $a_n\downarrow0$, $A_n=a_n^{-2}L_n$, wird
\emph{postuliert, nicht bewiesen}: Die Dirichlet-Formen $\mathcal E_n$
konvergieren im Mosco-Sinn gegen
$\mathcal E[f]=\int\nabla f^{T}C(x)\nabla f\,d^3x$ mit $C(x)\approx c^2\mathbb1$.

\subsection{Der bedingte hyperbolische Limes}

Unter fünf unbewiesenen Annahmen (Mosco-Konvergenz, gleichmäßige
Energiestabilität, konvergente Anfangsdaten, verschwindende
Konsistenzfehler, Kompaktheit) ergibt sich im schwachen Grenzwert
$\partial_t^2u-c^2\Delta u=0$ — ,,keiner der fünf Punkte ist in diesem
Preprint tatsächlich gezeigt.``

% ============================================================
\section{Die alternative Konstruktion — Kontinuum als kanonischer Pullback statt als Grenzwert einer Folge}

Anstatt eine Folge diskreter Approximationen mit schrumpfender
Gitterkonstante zu betrachten, wird die kompakte Struktur von vornherein
als fundamental behandelt. Der Übergang zur nichtkompakten Beschreibung
erfolgt über einen \textbf{kanonischen Funktionen-Pullback} entlang einer
festen Überlagerungsprojektion $p:\mathbb{R}\to S^1$.

\begin{quellbox}[title={Der Pullback als exakte Abbildung}]
$\Phi=p^{*}:C(S^1)\to C(\mathbb{R})$ ist eine exakte,
unendlich-dimensionale, stetige, algebrenerhaltende Abbildung — keine
Familie von Approximationen. Jede periodische Funktion auf $S^1$ liefert
sofort eine periodische Funktion auf $\mathbb{R}$; die verlustbehaftete
Richtung liegt bei $\mathbb{R}\to S^1$, nicht bei der Kompaktifizierung
selbst.
\end{quellbox}

\subsection{Warum dies methodisch etwas anderes ist als die Kategorienfehler-Blockade}

Die Kategorienfehler-Blockade richtet sich gegen die Identifikation zweier
Kettenräume \emph{unterschiedlichen Grades} innerhalb eines diskreten
Kettenkomplexes. Der kanonische Funktionen-Pullback fällt strukturell
nicht in diese Kategorie: Er identifiziert nicht zwei verschiedene
Kettengrade, sondern zwei \emph{Darstellungen derselben
Funktionenalgebra}. Die allgemeine methodische Sorge bleibt für echte
Kettengrad-Wechsel berechtigt; sie beschreibt schlicht nicht den Fall
eines expliziten, deklarierten, operationserhaltenden Pullbacks zwischen
Funktionenräumen.

\begin{kernbox}[title={Präzisierung: zwei verschiedene Bedeutungen von ,,Kontinuum``}]
Der Pullback $\Phi=p^*:L^2(S^1_m)\to\mathrm{Per}_{R_m}(\mathbb{R}_t)$,
$(\Phi f)(t)=f(t\bmod R_m)$, liefert ein \emph{Kontinuum der
Koordinatendarstellung}: Ein hinreichend regulärer Repräsentant von $f$
ist an jedem Punkt $t\in\mathbb{R}$ statt nur modulo $R_m$ auswertbar,
mit nicht-periodischen Testfunktionen vergleichbar, und (fast überall)
differenzierbar. Da es sich hier durchweg um Funktionen \emph{einer}
reellen Variablen $t$ handelt (der Fall $n=1$), sichert der
Sobolev-Einbettungssatz bereits für $f\in H^1$ die Einbettung
$H^1(I)\hookrightarrow C^{0,1/2}(I)$ (Hölder-$1/2$-stetig) auf jedem
beschränkten Intervall $I$ bzw.\ auf $S^1_m$ selbst: Ein solcher
Repräsentant ist stetig, sogar absolut stetig, und damit nach dem
Hauptsatz der Differential- und Integralrechnung für Lebesgue-Integrale
\emph{fast überall} klassisch differenzierbar, mit Ableitung gleich der
schwachen Ableitung fast überall (Beispiel: $f(t)=|t|$ liegt in $H^1$,
ist nicht $C^1$, aber überall stetig und außer bei einem Punkt klassisch
differenzierbar). Volle $C^1$-Regularität ist eine strengere, separate
Bedingung -- sie verlangt eine \emph{stetige} (nicht nur fast überall
definierte) Ableitung, keine zusätzliche Differenzierbarkeit an sich.
Für höhere Dimensionen $n\geq2$ gilt diese Einbettung nicht mehr
($H^1$-Funktionen können dort unbeschränkt sein); dieser Fall liegt hier
jedoch nicht vor. Diese Eigenschaft gehört dem Repräsentanten, nicht dem
bloßen $L^2$-Element selbst: $L^2(S^1_m)$ (ebenso $\mathrm{Per}_{R_m}(\mathbb{R}_t)$)
besteht aus Äquivalenzklassen fast-überall-gleicher Funktionen, für die
ein Funktionswert an einem festen Punkt $t$ ohne zusätzliche Regularität
schlicht nicht wohldefiniert ist -- unabhängig davon, ob man auf $S^1_m$
bleibt oder den Pullback nach $\mathbb{R}_t$ vollzieht. Er liefert
\emph{kein} Kontinuum des Zustandsraums: Da $\Phi f$ $R_m$-periodisch
ist, gilt $\int_{\mathbb{R}}|\Phi f|^2\,dt=\infty$ für jedes $f\neq0$ --
$\Phi f$ liegt nicht in gewöhnlichem $L^2(\mathbb{R}_t)$ mit Lebesgue-Maß
über die ganze Linie. Der korrekte Zielraum $\mathrm{Per}_{R_m}(\mathbb{R}_t)$
(periodischer $L^2_\mathrm{loc}$, Norm über eine Periode) trägt dieselbe
abzählbare Fourier-Basis $\{e_n\}_{n\in\mathbb{Z}}$ und dieselbe
Parseval-Identität wie $L^2(S^1_m)$ selbst -- die Isometrie ist ein
Darstellungswechsel innerhalb \emph{isomorpher} Hilberträume, keine
Erweiterung zu einem größeren Zustandsraum. Beide Aussagen -- Kontinuum
der Koordinaten, kein Kontinuum
des Zustandsraums -- sind gleichzeitig wahr und widersprechen sich nicht;
sie beantworten nur verschiedene Fragen. Diese Unterscheidung schärft,
nicht widerlegt, die im vorigen Absatz beschriebene Eigenschaft des
Pullbacks als reinem Darstellungswechsel.
\end{kernbox}

\subsection{Der Unterschied zur Mosco-Konvergenz-Strategie}

Der wesentliche Unterschied liegt genau in den Bausteinen ,,Mosco-Kon\-vergenz``
und ,,bedingter hyperbolischer Limes``: Dort ist eine skalierte
Verfeinerungsfolge mit einer \emph{postulierten, nicht bewiesenen}
Mosco-Konvergenz notwendig. Bei der Pullback-Konstruktion entfällt diese
Verfeinerungsmaschinerie vollständig — es gibt keine Folge $a_n\downarrow0$,
keinen reskalierten Operator $A_n$, keine Mosco-Verpflichtung. Der Preis:
Die kompakte Struktur selbst wird nicht aus etwas Fundamentalerem
hergeleitet, sondern gesetzt.

\subsection{Was die Trivialität der spektralen Dimension für einen Operator mit unendlichem diskreten Spektrum bedeutet}

Der Trivialitätssatz betrifft ausdrücklich einen \emph{endlichen} Operator
auf \emph{festem endlichem Träger}. Ein Operator mit von vornherein
unendlich vielen diskreten Eigenwerten auf einer kompakten Mannigfaltigkeit
— etwa der Laplace-Operator auf $S^1$ mit Spektrum $\{k^2:k\in\mathbb{Z}\}$
— erfüllt diese Voraussetzung nicht: $N=\dim\mathcal H=\infty$, das
Spektrum divergiert unbeschränkt statt bei einem endlichen Größtwert
abzubrechen. Bereits die Taylor-Entwicklung $P(\sigma)\approx1-\sigma\bar\lambda$
bricht zusammen, da ein Mittelwert über unendlich viele, unbeschränkt
wachsende Eigenwerte nicht wohldefiniert ist. Der Trivialitätssatz trifft
auf einen solchen Operator daher schlicht nicht zu — weder bestätigt noch
widerlegt, seine Voraussetzung liegt nicht vor.

% ============================================================
\section{Zusammenfassung}

\begin{tabular}{p{4.2cm}p{5.7cm}p{5.7cm}}
\toprule
& \textbf{Diskrete Kettenkomplex-Strategie} & \textbf{Kanonische Pullback-Strategie}\\
\midrule
Ausgangspunkt & endlicher Zellkomplex & kompakte Struktur als fundamental gesetzt\\
Kategorienfehler & durch explizites Lemma blockiert & Frage stellt sich nicht (Darstellungswechsel, kein Gradwechsel)\\
Geometrisches Scheitern & präzise lokalisiert (algebraisch bleibt gültig) & entfällt (keine Einbettungsfrage)\\
Kontinuumsübergang & Verfeinerungsfolge, Mosco-Konvergenz \emph{postuliert} & einzelner exakter Pullback, keine Folge nötig\\
Offene Voraussetzungen & fünf unbewiesene Annahmen für hyperbolischen Limes & Kompaktheit der Grundstruktur selbst ungeklärt\\
\bottomrule
\end{tabular}

\medskip
Beide Strategien haben ihren Preis: Die erste zahlt mit unbewiesenen
Konvergenzannahmen für den Kontinuumslimes, die zweite zahlt damit, dass
die Kompaktheit der fundamentalen Struktur selbst als Ausgangssetzung
steht. Welcher Preis der akzeptablere ist, hängt vom jeweiligen Anspruch
der Konstruktion ab.

% ============================================================
\section{Literatur}

\begin{thebibliography}{9}
\bibitem{Krueger2026}
Krüger, M. (2026). \emph{Dimension-Graded Geometry and Typed Emergence on
Aperiodic Carriers: From Cellular Degree, Holonomy, and Regge Curvature
Support to a Falsifiable Spacetime-Limit Programme in an HLV Testbed.}
Zenodo. DOI: 10.5281/zenodo.22105698.

\bibitem{Regge1961}
Regge, T. (1961). General relativity without coordinates.
\emph{Il Nuovo Cimento}, 19(3), 558--571.

\bibitem{Wilson1974}
Wilson, K. G. (1974). Confinement of quarks.
\emph{Physical Review D}, 10, 2445--2459. — Ursprungsarbeit der
Gitter-Eichtheorie; Grundlage für den Kanten-Transport-/Flächen-Holonomie-Baustein
(Abschnitt 1.3).

\bibitem{BahrDittrich2009}
Bahr, B., \& Dittrich, B. (2009). Regge calculus from a new angle.
\emph{New Journal of Physics}, 12, 033010. arXiv:0907.4325. — Verallgemeinerter
Regge-Zugang mit konstant gekrümmten statt flachen Simplizes; siehe die
Einordnung in Abschnitt 1.4.

\bibitem{DittrichHoehn2012}
Dittrich, B., \& Höhn, P. A. (2012). From covariant to canonical formulations
of discrete gravity. \emph{Classical and Quantum Gravity}, 27, 155001.
arXiv:0912.1817. — Verallgemeinerter Regge-Zugang mit variabler
Triangulierung über die Zeit; siehe Abschnitt 1.4.

\bibitem{Hohn2015}
Höhn, P. A. (2015). Canonical linearized Regge calculus: counting lattice
gravitons with Pachner moves. \emph{Physical Review D}, 91, 124034.

\bibitem{DesbrunEtAl2005}
Desbrun, M., Hirani, A. N., Leok, M., \& Marsden, J. E. (2005). Discrete
exterior calculus. arXiv:math/0508341. — Grundlagenarbeit für die
diskrete Außenrechnung und die Hodge-Positivität aus Abschnitt 1.6.

\bibitem{FriedmanSorkin1980}
Friedman, J. L., \& Sorkin, R. D. (1980). Spin $1/2$ from gravity.
\emph{Physical Review Letters}, 44, 1100.

\bibitem{AmbjornJurkiewiczLoll2005}
Ambjørn, J., Jurkiewicz, J., \& Loll, R. (2005). Spectral dimension of the
universe. \emph{Physical Review Letters}, 95, 171301. arXiv:hep-th/0505113.
— Empirische Illustration des Trivialitätssatzes aus Abschnitt 1.7:
Messungen an endlichen kausalen dynamischen Triangulierungen zeigen genau
das dort beschriebene Zwischenskalen-Plateau samt endlichkeitsbedingtem
Abfall auf null bei kleinem und großem Diffusionsparameter.

\bibitem{AmbjornJurkiewiczLollReview}
Ambjørn, J., Görlich, A., Jurkiewicz, J., \& Loll, R. (2012). Nonperturbative
quantum gravity. \emph{Physics Reports}, 519(4--5), 127--210.
arXiv:1203.3591. — Übersichtsartikel zu kausalen dynamischen
Triangulierungen mit ausführlicher Diskussion endlichkeitsbedingter
Effekte auf die spektrale Dimension.

\bibitem{MoscoConvergence1994}
Mosco, U. (1994). Composite media and asymptotic Dirichlet forms.
\emph{Journal of Functional Analysis}, 123(2), 368--421. — Ursprungsarbeit
zum in Abschnitt 1.8 verwendeten Konvergenzbegriff.

\bibitem{DiscreteToContinuumMosco2025}
Beitrag zur Mosco-Konvergenz diskreter Graph-Dirichlet-Formen auf dem
flachen Torus $\mathbb{T}^n$ gegen ein kontinuierliches Zielfunktional
(arXiv:2511.16486, 2025). — Behandelt exakt den in Abschnitt 1.8
beschriebenen Konstruktionstyp (schrumpfende Gitterkonstante $\varepsilon\to0$
auf einem torusförmigen Gitter) und zeigt, unter welchen zusätzlichen
technischen Bedingungen eine solche Mosco-Konvergenz tatsächlich beweisbar
ist — eine nützliche Gegenprobe zur hier referierten, unbewiesenen
Verpflichtung.

\bibitem{KaluzaKlein1926}
Klein, O. (1926). Quantentheorie und fünfdimensionale Relativitätstheorie.
\emph{Zeitschrift für Physik}, 37(12), 895--906. — Klassisches Beispiel
einer Pullback-artigen Konstruktion in der Physik: die Fourier-Zerlegung
eines Feldes entlang einer kompakten Kreisfaser, strukturell verwandt mit
dem in Abschnitt 2 diskutierten kanonischen Funktionen-Pullback.
\end{thebibliography}
